\chapter{Descripción del problema}\label{cap:descripcion}

La computación en la nube ha transformado la forma en que se despliega y gestiona
la infraestructura tecnológica, ofreciendo acceso bajo demanda a recursos de cómputo,
almacenamiento y red sin necesidad de inversión en hardware propio. Sin embargo,
detrás de esta promesa de flexibilidad subyace un problema estructural que condiciona
profundamente las decisiones de quienes adoptan estos servicios: el \textbf{cloud
lock-in}. Este capítulo analiza dicho problema en su dimensión completa, presenta
la solución propuesta y delimita el alcance de este trabajo.

\section{Contexto y problemática}

La adopción de servicios cloud ha crecido de forma sostenida durante la última
década. Según el informe anual de Flexera sobre el estado de la nube
\cite{flexera2026}, el 73\% de las organizaciones encuestadas operan en entornos
híbridos o multi-cloud, lo que refleja tanto la madurez del mercado como la
creciente conciencia sobre los riesgos de depender de un único proveedor. Sin
embargo, este dato contrasta con la dificultad real que supone mover cargas de
trabajo entre plataformas una vez que se ha iniciado el despliegue.

El \textbf{cloud lock-in}, o dependencia del proveedor cloud, se produce cuando
una organización o individuo queda vinculado a una plataforma concreta de forma
que migrar a otra resulta técnicamente complejo, económicamente costoso o
ambas cosas a la vez. Esta dependencia no es accidental: los proveedores ofrecen
servicios gestionados, integraciones nativas y modelos de precios que incentivan
el uso de sus propias herramientas, generando un ecosistema difícil de abandonar.

En este capítulo, el problema se estructura en tres dimensiones que se refuerzan mutuamente:

\begin{itemize}
    \item \textbf{Lock-in de infraestructura:} los recursos desplegados en un
    proveedor —instancias de cómputo, bases de datos gestionadas, configuraciones
    de red— utilizan servicios, nomenclaturas y modelos de configuración específicos
    de esa plataforma. Migrar una arquitectura de \textbf{AWS} a
    \textbf{Microsoft Azure}, por ejemplo, no es una cuestión de exportar y
    reimportar configuraciones: requiere rediseñar los recursos en los términos
    propios del proveedor destino.

    \item \textbf{Lock-in de conocimiento:} aunque los conceptos fundamentales
    de la computación en la nube son transferibles entre plataformas, las
    tecnologías concretas no lo son. Un profesional formado en AWS —familiarizado
    con \textbf{EC2}, \textbf{RDS} o \textbf{VPC}— debe afrontar un proceso de
    reaprendizaje significativo para operar con los servicios equivalentes de Azure
    o \textbf{Google Cloud Platform (GCP)}. Este coste de conocimiento desincentiva
    la exploración de alternativas y consolida la dependencia.

    \item \textbf{Lock-in económico:} la estimación de costes cloud es opaca y
    compleja. Aunque los proveedores ofrecen calculadoras de precios, su uso
    requiere conocimiento previo de los servicios a contratar y criterio suficiente
    para modelar correctamente el escenario de despliegue. Comparar el coste de
    una misma arquitectura en distintos proveedores exige dominar simultáneamente
    los modelos de precios de cada uno, tarea que en la práctica recae en
    especialistas o empresas externas \cite{flexera2026}.
\end{itemize}

A estas tres dimensiones se añade la \textbf{complejidad de seguridad} inherente
a la nube. Todos los proveedores cloud operan bajo un modelo de
\hyperlink{shared-responsibility}{\textbf{responsabilidad compartida}}: el proveedor garantiza la seguridad de la
infraestructura subyacente, pero la seguridad de lo que el usuario despliega sobre
ella —configuración de permisos, exposición de puertos, cifrado de datos— es
responsabilidad exclusiva de quien gestiona los recursos \cite{awsSharedResponsibility}.
Este modelo, combinado con la curva de aprendizaje de cada plataforma, genera un
riesgo adicional para equipos con experiencia limitada en la nube.

\section{Magnitud del problema}

El impacto del cloud lock-in no se limita a la dificultad técnica de una migración
puntual. Sus consecuencias afectan a la capacidad de negociación, a la eficiencia
operativa y a la sostenibilidad económica de proyectos de cualquier escala.

En el ámbito económico, el sobrecoste por uso ineficiente o no planificado de
recursos cloud es un problema extendido. El mismo informe de Flexera
\cite{flexera2026} estima que las organizaciones desperdician de media un 29\%
de su gasto cloud, en gran parte por falta de visibilidad y herramientas de
control adecuadas. La ausencia de mecanismos que permitan comparar costes entre
proveedores antes de comprometerse con uno agrava esta situación: la elección
del proveedor se realiza frecuentemente por familiaridad o inercia, no por
criterio económico fundamentado.

En el ámbito organizacional, la dependencia de un proveedor limita la capacidad
de respuesta ante cambios de precios, interrupciones de servicio o evoluciones
del mercado. Empresas que han centralizado su infraestructura en un único cloud
se encuentran en una posición negociadora débil y asumen el riesgo de que
decisiones unilaterales del proveedor —cambios de tarifas, discontinuación de
servicios— afecten directamente a su operativa.

\section{Solución propuesta}

Frente a este problema, se propone \textbf{CloudPorter}: una herramienta de línea de
comandos que permite describir infraestructura cloud de forma agnóstica al
proveedor y desplegarla en distintas plataformas a partir de una única definición.
El objetivo último es dotar de interoperabilidad a la infraestructura cloud,
la capacidad de trasladar una misma definición de un proveedor a otro sin reescribirla:
una propiedad de la que el ecosistema actual carece y que constituye la raíz del problema
del lock-in.

El núcleo de CloudPorter es el \hyperlink{cloudporter-manifest}{\textbf{CloudPorter Manifest}}, un fichero en formato
\textbf{YAML} que describe los recursos de infraestructura en términos genéricos, independientes de cualquier proveedor. A partir
de este manifiesto, CloudPorter traduce la definición a las plantillas de infraestructura
como código propias de cada plataforma, abstrayendo al usuario de los detalles
específicos de cada cloud.

Además de la traducción y el despliegue, CloudPorter incorpora la posibilidad de estimación
de costes que permite comparar el coste aproximado de una arquitectura en distintos
proveedores antes de ejecutar el despliegue. De este modo, la elección del proveedor
puede fundamentarse en criterios económicos objetivos, reduciendo la opacidad que
caracteriza hoy a este proceso.

\section{Alcance y restricciones}

El problema del cloud lock-in es amplio y sus dimensiones abarcan aspectos que van
más allá de lo abordable en el marco de un único trabajo. Por ello, se establece
explícitamente qué entra en el alcance de este TFG y qué queda como trabajo futuro.

En el \textbf{alcance del trabajo} se incluye:

\begin{itemize}
    \item \textbf{Soporte a dos proveedores cloud:} AWS y Microsoft Azure como
    plataformas de despliegue objetivo en la versión inicial de CloudPorter.

    \item \textbf{Recursos de infraestructura básicos:} cómputo, configuración
    de red —generada automáticamente, sin requerir declaración explícita del
    usuario— y bases de datos relacionales.

    \item \textbf{Estimación comparativa de costes:} módulo que permite comparar
    el coste aproximado de un despliegue entre los proveedores soportados.

    \item \textbf{Interfaz de línea de comandos:} como interfaz principal de
    interacción con la herramienta.
\end{itemize}

Quedan \textbf{fuera del alcance} de este trabajo:

\begin{itemize}
    \item \textbf{Soporte a Google Cloud Platform} y otros proveedores: el soporte
    a dos clouds es suficiente para demostrar la viabilidad del enfoque agnóstico.

    \item \textbf{Almacenamiento de objetos:} los servicios de almacenamiento,
    como \textbf{AWS S3} o \textbf{Azure Blob Storage}, implican lock-in a nivel
    de código, no solo de infraestructura, lo que excede el alcance de este trabajo.

    \item \textbf{Análisis de seguridad avanzado:} aunque se reconoce la
    complejidad del modelo de responsabilidad compartida, su abordaje requiere
    un modelado de políticas específico por proveedor que excede el alcance de
    este trabajo.

    \item \textbf{Interfaz web o plataforma:} la CLI es suficiente para validar
    el núcleo de la herramienta; una interfaz gráfica no aporta evidencia adicional
    sobre la viabilidad del enfoque.
\end{itemize}
